%\newcommand{\CLASSINPUTrighttextmargin}{0.62cm}
\documentclass[conference]{IEEEtran}
\IEEEoverridecommandlockouts
\usepackage{cite}
\usepackage{float}
\usepackage{amsmath,amssymb,amsfonts}
\usepackage{algorithmic}
\usepackage[boxed,lined,ruled]{algorithm2e}
\usepackage{algpseudocode}
\usepackage{graphicx}
\usepackage{textcomp}
\usepackage{xcolor}
\usepackage{multirow}
\usepackage{scalefnt}
\usepackage{booktabs}
\usepackage[brazil]{babel}

\begin{document}

\title{Processamento Digital de Sinais: Nome do Algoritmo Utilizado}

\author{Nome do Autor 1, Nome do Autor 2, Nome do Autor 3 e Nome do Autor 4}

\maketitle

\section{Fundamentação Teórica}

Explicar o algoritmo de processamento digital de sinais utilizado. 

Citar as referências bibliográficas utilizadas para a compreensão da fundamentação teórica do algoritmo. Por exemplo, para citar uma referência que trata de aplicações da \textit{Fast Fourier Transform} (FFT), você deve cadastrar os dados bibliográficos no arquivo ref.bib e chamar a citação da seguinte forma~\cite{cooley1969}. 

Detalhar e explicar seu equacionamento. Criar um pseudo-código (use o pacote algorithm) ou um fluxograma (usar uma figura) do algoritmo. A Figura~\ref{fig:flowchart} traz um exemplo de fluxograma.

\begin{figure}[h]
\centering
\includegraphics[width=0.45\textwidth]{ex_fig1.pdf}
\caption{Exemplo de Fluxograma.}
\label{fig:flowchart}
\end{figure}

\textcolor{red}{Obs.: Será verificado o uso de IA Generativa. Somente será aceito para fins de correção ortográfica/gramatical. Escrita feita por IA Generativa não é permitida.}

\textcolor{blue}{Obs.: Esse documento deve ter entre 5 e 8 páginas, sem contar o Apêncice.}

\section{Análise dos Sinais}

Apresente o(s) sinal(is) que seu grupo está utilizando. Utilize figuras com boa resolução (evitar .png -- busque usar formatos vetorizados). A Figura~\ref{fig:grafico} traz um exemplo de gráfico de sinal.

\begin{figure}[h]
\centering
\includegraphics[width=0.45\textwidth]{ex_fig2.pdf}
\caption{Exemplo de Gráfico para Sinais.}
\label{fig:grafico}
\end{figure}

Explique as parametrizações testadas para o algoritmo de processamento digital de sinais utilizado (se houver). 

Demonstre, de preferência graficamente, as saídas do algoritmo de processamento digital de sinais utilizado. Ou traga valores tabelados. A Tabela~\ref{tab:results} exemplifica a tabulação de dados para apresentação de resultados.

\begin{table}[!h]
\renewcommand{\arraystretch}{1.2}
\caption{Dados tabulados.}
\label{tab:results}
\centering
\begin{tabular}{|c|ccccc|}
\hline
\multirow{2}{*}{\textbf{Zone}} & \multicolumn{5}{c|}{\textbf{MRE}} \\
 & \textbf{6-pulse} & \textbf{12-pulse} & \textbf{SFC} & \textbf{DC motor} & \textbf{TCR} \\
\hline
\textbf{1} & 15.2\% & 14.4\% & 32.2\% & 5.9\% & 35.2\% \\
\textbf{2} & 10.3\% & 12.0\%  & 7.1\% & 5.3\%  & 19.7\% \\
\textbf{3} & 0.6\% & 2.3\%  & 2.2\% & 1.3\%  & 0.5\% \\
\hline
\textbf{General} & 6.9\%  & 8.2\%  & 9.4\% & 3.7\%  & 14\% \\
\hline
\end{tabular}
\end{table}

Apresente e discuta os resultados finais de seu projeto. Pode usar gráficos e/ou tabelas. Entretanto, precisam ser devidamente explicados.

\section{Conclusões}

Apresente as conclusões obtidas pelo grupo em relação à aplicação do algoritmo. 

\bibliographystyle{IEEEtran}
\bibliography{ref}

\appendix

Trazer o código da função da técnica de processamento digital de sinais implementada. Se o grupo implementou com base em algum repositório do GitHub, colocar tanto o código adaptado quanto o link desse repositório. Se o grupo foi auxiliado por IA Generativa para implementar o código, colocar o prompt usado e qual ferramenta foi usada (Ex.: Claude, Gemini, ChatGPT).

\end{document}
